%-----------------------------------------------------------------------------%
\chapter{PERANCANGAN}
\label{bab:4}
%-----------------------------------------------------------------------------%

\section{Perancangan URI}\label{perancangan-uri}

Setiap entitas yang berbeda perlu diberikan URI yang berbeda. Penggabungan KG sangat umum dilakukan
karena secara struktur dapat dilakukan dengan mudah, dan dapat memperkuat data dari masing-masing
KG. KG memiliki entitas-entitas dengan URI yang unik, oleh karena itu saat menggabungkan sebuah KG A
dan sebuah KG B, entitas yang berbeda harus memiliki URI yang berbeda. Sebuah masalah dapat terjadi
apabila URI tidak dirancang dengan benar. Sebagai contoh, apabila terdapat entitas dari KG A dengan
URI \texttt{http://example.org/trump} dengan maksud Donald Trump mantan presiden Amerika Serikat,
dan entitas dari KG B dengan URI yang sama dengan maksud permainan kartu Trump, saat kedua KG ini
digabung akan menjadi informasi yang berkontradiksi.

Pada penelitian ini, semua URI menggunakan prefix \texttt{http://example.org/legal/} yang
selanjutnya akan disebut \emph{basePrefix}. Prefix ini dapat diganti dalam program agar dapat
disesuaikan oleh pengguna program. Semua entitas dokumen dan komponennya menggunakan prefix URI
\texttt{http://example.org/legal/peraturan} yang selanjutnya disebut \emph{peraturanPrefix}, dan
semua resource lainnya menggunakan \texttt{http://example.org/legal/ontology} yang selanjutnya
disebut \emph{ontoPrefix}.

\section{Perancangan URI Class}\label{perancangan-uri-class}

URI Class adalah URI yang menjelaskan \emph{class} dari entitas tersebut. Semua URI Class mengikuti
pola \texttt{\{ontoPrefix\}/\{className\}} dimana \texttt{className} adalah nama \emph{class}
tersebut. Berikut adalah semua \emph{class} komponen beserta URI dan deskripsinya.

\begin{longtable}[c]{@{}ll@{}} \toprule\addlinespace Nama Class & Deskripsi
  \\\addlinespace \midrule\endhead \texttt{Peraturan}      & Peraturan perundang-undangan
  \\\addlinespace \texttt{DaftarBab}      & Daftar satu atau lebih bab \\\addlinespace \texttt{Bab}
                                                 & Bab                                                                  \\\addlinespace \texttt{DaftarBagian}   & Daftar satu atau lebih bagian \\\addlinespace
  \texttt{Bagian}                                & Bagian                                                               \\\addlinespace \texttt{DaftarParagraf} & Daftar satu atau lebih
  paragraf                                                                                                              \\\addlinespace \texttt{Paragraf}       & Paragraf \\\addlinespace \texttt{DaftarPasal}
                                                 & Daftar satu atau lebih pasal                                         \\\addlinespace \texttt{Pasal}          & Pasal \\\addlinespace
  \texttt{VersiPasal}                            & Versi dari sebuah pasal. Sebuah pasal dapat memiliki satu atau lebih
  versi dari hasil amandemen                                                                                            \\\addlinespace \texttt{DaftarAyat}     & Daftar satu atau lebih ayat
  \\\addlinespace \texttt{Ayat}           & Ayat \\\addlinespace \texttt{DaftarHuruf}    & Daftar
  satu atau lebih huruf atau nomor                                                                                      \\\addlinespace \texttt{Huruf}          & Huruf atau nomor
  \\\addlinespace \texttt{Menimbang}      & Hal yang ditimbang oleh suatu peraturan \\\addlinespace
  \texttt{Mengingat}                             & Hal yang diingat oleh suatu peraturan                                \\\addlinespace \texttt{Segmen}
                                                 & Segmen teks. Dapat memiliki rujukan ke suatu komponen peraturan      \\\addlinespace
  \bottomrule
\end{longtable}

\section{Perancangan URI Komponen Peraturan
  Perundang-undangan}\label{perancangan-uri-komponen-peraturan-perundang-undangan}

URI peraturan perundang-undangan diawali oleh \emph{peraturanPrefix}, kemudian diikuti oleh string
yang mengidentifikasi peraturan tersebut. URI peraturan ini akan digunakan oleh entitas sebagai
prefix yang selanjutnya akan disebut \emph{docURI}. Berikut adalah jenis peraturan
perundang-undangan yang diimplementasi pada penelitian ini beserta pola URI nya. Variabel akan
dikurung dengan kurung kurawal \texttt{\{\}}.

\begin{longtable}[c]{@{}lll@{}} \toprule\addlinespace Jenis Peraturan
                                                      & Pola URI                                                                       & Penjelasan variabel             \\\addlinespace
  \midrule\endhead Undang-Undang Dasar 1945           &
  \texttt{\{peraturanPrefix\}/uud}                    &                                                                                                                  \\\addlinespace Undang-Undang
                                                      & \texttt{\{peraturanPrefix\}/uu/\{tahun\}/\{nomor\}}                            & \texttt{tahun}: tahun peraturan
  \\\addlinespace &                                                     & \texttt{nomor}: nomor
  peraturan                                                                                                                                                              \\\addlinespace Peraturan Pemerintah                                                  &
  \texttt{\{peraturanPrefix\}/pp/\{tahun\}/\{nomor\}} & \texttt{tahun}: tahun
  peraturan                                                                                                                                                              \\\addlinespace &                                                     & \texttt{nomor}:
  nomor peraturan                                                                                                                                                        \\\addlinespace Peraturan Daerah Provinsi DKI Jakarta
                                                      & \texttt{\{peraturanPrefix\}/perda\_provinsi\_dki\_jakarta/\{tahun\}/\{nomor\}} & \texttt{tahun}:
  tahun peraturan                                                                                                                                                        \\\addlinespace &                                                     &
  \texttt{nomor}: nomor peraturan                                                                                                                                        \\\addlinespace Peraturan Gubernur DKI Jakarta
                                                      & \texttt{\{peraturanPrefix\}/pergub\_dki\_jakarta/\{tahun\}/\{nomor\}}          & \texttt{tahun}: tahun
  peraturan                                                                                                                                                              \\\addlinespace &                                                     & \texttt{nomor}:
  nomor peraturan                                                                                                                                                        \\\addlinespace Peraturan Walikota Malang
                                                      & \texttt{\{peraturanPrefix\}/perwali\_malang/\{tahun\}/\{nomor\}}               & \texttt{tahun}: tahun
  peraturan                                                                                                                                                              \\\addlinespace &                                                     & \texttt{nomor}:
  nomor peraturan                                                                                                                                                        \\\addlinespace
  \bottomrule
\end{longtable}

Setiap komponen dokumen merupakan entitas, dan memiliki URI. URI sebuah komponen didahului oleh
\emph{docURI} dan semua komponen yang mengatasinya. Berikut adalah jenis komponen peraturan
perundang-undangan beserta pola URI nya.

\begin{longtable}[c]{@{}lll@{}} \toprule\addlinespace Class Komponen & Pola URI
                                                      & Penjelasan variabel                                        \\\addlinespace \midrule\endhead \texttt{DaftarBab}      &
  \texttt{\{docURI\}/bab}                             & \texttt{docURI}: URI peraturan yang mengatasi
  \\\addlinespace \texttt{Bab}            & \texttt{\{daftarBabURI\}/\{no\}}      &
  \texttt{daftarBabURI}: URI \texttt{DaftarBab} yang mengatasi                                                     \\\addlinespace &
                                                      & \texttt{no}: Nomor bab                                     \\\addlinespace \texttt{DaftarBagian}   & \texttt{\{babURI\}/bagian}
                                                      & \texttt{babURI}: URI \texttt{Bab} yang mengatasi           \\\addlinespace \texttt{Bagian}         &
  \texttt{\{daftarBagianURI\}/\{no\}}                 & \texttt{daftarBagianURI}: URI \texttt{DaftarBagian} yang
  mengatasi                                                                                                        \\\addlinespace &                                       & \texttt{no}: Nomor bagian
  \\\addlinespace \texttt{DaftarParagraf} & \texttt{\{bagianURI\}/paragraf}       &
  \texttt{bagianURI}: URI \texttt{Bagian} yang mengatasi                                                           \\\addlinespace \texttt{Paragraf}       &
  \texttt{\{daftarParagrafURI\}/\{no\}}               & \texttt{daftarParagrafURI}: URI \texttt{DaftarBagian} yang
  mengatasi                                                                                                        \\\addlinespace &                                       & \texttt{no}: Nomor paragraf
  \\\addlinespace \texttt{DaftarPasal}    & \texttt{\{parentURI\}/daftarpasal}    &
  \texttt{parentURI}: URI \texttt{Bab}, \texttt{Bagian}, \texttt{Paragraf} yang mengatasi
  \\\addlinespace \texttt{Pasal}          & \texttt{\{docURI\}/pasal/\{no\}}      & \texttt{docURI}:
  URI peraturan yang mengatasi                                                                                     \\\addlinespace &                                       &
  \texttt{no}: Nomor pasal                                                                                         \\\addlinespace \texttt{VersiPasal}     &
  \texttt{\{pasalURI\}/versi/\{date\}}                & \texttt{pasalURI}: URI \texttt{Pasal} yang mengatasi
  \\\addlinespace &                                       & \texttt{date}: Tanggal disahkan
  \\\addlinespace \texttt{DaftarAyat}     & \texttt{\{pasalURI\}/ayat}            &
  \texttt{pasalURI}: URI \texttt{Pasal} yang mengatasi                                                             \\\addlinespace \texttt{Ayat}           &
  \texttt{\{daftarAyatURI\}/\{no\}}                   & \texttt{daftarAyatURI}: URI \texttt{DaftarAyat} yang
  mengatasi                                                                                                        \\\addlinespace &                                       & \texttt{no}: Nomor ayat
  \\\addlinespace \texttt{DaftarHuruf}    & \texttt{\{parentURI\}/huruf}          &
  \texttt{parentURI}: URI \texttt{Ayat}, \texttt{Huruf}, \texttt{VersiPasal}, \texttt{Menimbang},
  \texttt{Mengingat}, yang mengatasi                                                                               \\\addlinespace \texttt{Huruf}          &
  \texttt{\{daftarHurufURI\}/\{id\}}                  & \texttt{daftarHurufURI}: URI \texttt{DaftarHuruf} yang
  mengatasi                                                                                                        \\\addlinespace &                                       & \texttt{id}: Nomor atau huruf
  \\\addlinespace \texttt{Menimbang}      & \texttt{\{docURI\}/menimbang}         & \texttt{docURI}:
  URI peraturan yang mengatasi                                                                                     \\\addlinespace \texttt{Mengingat}      &
  \texttt{\{docURI\}/mengingat}                       & \texttt{docURI}: URI peraturan yang mengatasi
  \\\addlinespace \texttt{Segmen}         & \texttt{\{parentURI\}/\{name\}}       &
  \texttt{parentURI}: URI \texttt{Ayat}, \texttt{Huruf}, \texttt{DaftarHuruf}, \texttt{VersiPasal},
  \texttt{Menimbang}, \texttt{Mengingat}, yang mengatasi                                                           \\\addlinespace &
                                                      & \texttt{name}: Nama teks                                   \\\addlinespace
  \bottomrule
\end{longtable}

\section{Perancangan URI Properti}\label{perancangan-uri-properti}

Properti ditandai dengan URI yang memiliki pola \texttt{\{ontoPrefix\}/\{name\}} dimana
\texttt{name} adalah nama dari properti tersebut. Properti menghubungkan subjek dan objek. Subjek
adalah berupa URI komponen, sedangkan objek dapat berupa URI komponen, string, \emph{number}
(bilangan), atau \emph{date} (tanggal).

\begin{longtable}[c]{@{}llll@{}} \toprule\addlinespace Nama Properti
                                                                                    & Subjek                                                                                        & Objek                           & Deskripsi
  \\\addlinespace \midrule\endhead \texttt{nomor}                                & \texttt{Bab},
  \texttt{Bagian}, \texttt{Paragraf}, \texttt{Pasal}, \texttt{Ayat}, \texttt{Huruf} & string, number
                                                                                    & Nomor atau huruf \emph{identifier} dari komponen                                                                                               \\\addlinespace \texttt{teks}
                                                                                    & \texttt{Ayat}, \texttt{VersiPasal}, \texttt{Segmen}                                           & string
                                                                                    & teks dari komponen                                                                                                                             \\\addlinespace \texttt{bab}                                  &
  \texttt{DaftarBab}                                                                & \texttt{Bab}                                                                                  & Bab
  \\\addlinespace \texttt{bagian}                               & \texttt{DaftarBagian}
                                                                                    & \texttt{Bagian}                                                                               & Bagian                                         \\\addlinespace \texttt{paragraf}
                                                                                    & \texttt{DaftarParagraf}                                                                       & \texttt{Paragraf}               & Paragraf
  \\\addlinespace \texttt{pasal}                                & \texttt{DaftarPasal}
                                                                                    & \texttt{Pasal}                                                                                & Pasal                                          \\\addlinespace \texttt{ayat}
                                                                                    & \texttt{DaftarAyat}                                                                           & \texttt{Ayat}                   & Ayat
  \\\addlinespace \texttt{huruf}                                & \texttt{DaftarHuruf}
                                                                                    & \texttt{Huruf}                                                                                & Huruf                                          \\\addlinespace \texttt{segmen}
                                                                                    & \texttt{Mengigat},\texttt{Menimbang}, \texttt{Paragraf}, \texttt{VersiPasal}, \texttt{Huruf},
  \texttt{DaftarHuruf}                                                              & \texttt{Segmen}
                                                                                    & Segmen                                                                                                                                         \\\addlinespace \texttt{daftarBab}                            & \texttt{Peraturan}
                                                                                    & \texttt{DaftarBab}                                                                            & DaftarBab                                      \\\addlinespace \texttt{daftarBagian}
                                                                                    & \texttt{Bab}                                                                                  & \texttt{Daftarbagian}           & DaftarBagian
  \\\addlinespace \texttt{daftarParagraf}                       & \texttt{Bagian}
                                                                                    & \texttt{DaftarParagraf}                                                                       & DaftarParagraf                                 \\\addlinespace \texttt{daftarPasal}
                                                                                    & \texttt{Peraturan},\texttt{Bab}, \texttt{Paragraf}, \texttt{Bagian}                           &
  \texttt{DaftarPasal}                                                              & DaftarPasal                                                                                                                                    \\\addlinespace
  \texttt{daftarAyat}                                                               & \texttt{Pasal}
                                                                                    & \texttt{DaftarAyat}                                                                           & DaftarAyat                                     \\\addlinespace \texttt{daftarHuruf}
                                                                                    & \texttt{VersiPasal},\texttt{Ayat},\texttt{Huruf},\texttt{Menimbang},\texttt{Mengingat}        &
  \texttt{DaftarHuruf}                                                              & DaftarHuruf                                                                                                                                    \\\addlinespace \texttt{merujuk}
                                                                                    & \texttt{Segmen}                                                                               & entitas apapun                  & Teks merujuk
  suatu entitas                                                                                                                                                                                                                      \\\addlinespace \texttt{mengubah}                             & \texttt{Huruf}
                                                                                    & \texttt{VersiPasal}                                                                           & Pengubahan pasal                               \\\addlinespace \texttt{menghapus}
                                                                                    & \texttt{Huruf}                                                                                & \texttt{VersiPasal}             & Penghapusan
  pasal                                                                                                                                                                                                                              \\\addlinespace \texttt{menyisipkan}                          & \texttt{Huruf}
                                                                                    & \texttt{VersiPasal}                                                                           & Penyisipan pasal                               \\\addlinespace \texttt{tanggal}
                                                                                    & \texttt{VersiPasal}                                                                           & date                            & Tanggal
  \\\addlinespace \texttt{jenisVersi}                           & \texttt{VersiPasal}
                                                                                    & `orisinal', `penyisipan', `pengubahan', `penghapusan'                                         & Jenis versi
  \\\addlinespace \texttt{versi}                                & \texttt{Pasal}
                                                                                    & \texttt{VersiPasal}                                                                           & Versi suatu komponen                           \\\addlinespace \texttt{tentang}
                                                                                    & \texttt{Peraturan}                                                                            & string                          & Peraturan
  tentang                                                                                                                                                                                                                            \\\addlinespace \texttt{menimbang}                            & \texttt{Peraturan}
                                                                                    & \texttt{Menimbang}                                                                            & Peraturan menimbang                            \\\addlinespace \texttt{mengingat}
                                                                                    & \texttt{Peraturan}                                                                            & \texttt{Mengingat}              & Peraturan
  mengingat                                                                                                                                                                                                                          \\\addlinespace \texttt{disahkanPada}                         & \texttt{Peraturan}
                                                                                    & date                                                                                          & Peraturan disahkan pada tanggal                \\\addlinespace \texttt{disahkanDi}
                                                                                    & \texttt{Peraturan}                                                                            & string                          & Peraturan
  disahkan di lokasi                                                                                                                                                                                                                 \\\addlinespace \texttt{disahkanOleh}                         &
  \texttt{Peraturan}                                                                & string                                                                                        & Peraturan
  disahkan oleh                                                                                                                                                                                                                      \\\addlinespace \texttt{jabatanPengesah}                      & \texttt{Peraturan}
                                                                                    & string                                                                                        & Jabatan pengesah peraturan                     \\\addlinespace \texttt{bagianDari}
                                                                                    & komponen apapun                                                                               & komponen apapun                 & Komponen
  subjek adalah bagian dari komponen objek                                                                                                                                                                                           \\\addlinespace
  \bottomrule
\end{longtable}


